\documentclass[12pt]{article}

\usepackage[spanish]{babel}
\usepackage[utf8]{inputenc}
\usepackage[T1]{fontenc}
\usepackage{amsmath,amssymb}
\usepackage{geometry}
\usepackage{enumitem}
\usepackage{xcolor}
\usepackage{lmodern}

\geometry{letterpaper, margin=2cm}

\title{\textbf{Solucionario Guía 01}\\
	Fundamentos de Matemática Básica}
\author{}
\date{}

\begin{document}
	
	\maketitle
	
	%------------------------------------------------
	\section*{Parte I: Conjuntos numéricos}
	
	\begin{enumerate}
		
		\item 
		\[
		17+23+13+7=60 \quad \Rightarrow \quad \boxed{\mathbb{N}}
		\]
		
		\item 
		\[
		(12+18)+(8+2)=30+10=40 \quad \Rightarrow \quad \boxed{\mathbb{N}}
		\]
		
		\item 
		\[
		6(15+5)=6\cdot 20=120 \quad \Rightarrow \quad \boxed{\mathbb{N}}
		\]
		
		\item 
		\[
		8(20+10)-3(20+10)=240-90=150 \quad \Rightarrow \quad \boxed{\mathbb{N}}
		\]
		
		\item 
		\[
		5(14+6)+4(14+6)=100+80=180 \quad \Rightarrow \quad \boxed{\mathbb{N}}
		\]
		
		\item 
		\[
		7(12+8)+5(12+8)-3(12+8)=140+100-60=180 \quad \Rightarrow \quad \boxed{\mathbb{N}}
		\]
		
	\end{enumerate}
	
	%------------------------------------------------
	\section*{Parte II: Cierre de operaciones}
	
	\begin{enumerate}
		
		\item 
		\[
		8-13=-5 \notin \mathbb{N}
		\quad \Rightarrow \quad \boxed{\text{No está cerrada en } \mathbb{N}}
		\]
		
		\item 
		\[
		15 \div 5=3 \in \mathbb{Z}
		\quad \Rightarrow \quad \boxed{\text{Sí pertenece}}
		\]
		
		\item 
		\[
		\frac{2}{3}+\frac{5}{6}=\frac{4}{6}+\frac{5}{6}=\frac{9}{6}=\frac{3}{2}
		\quad \Rightarrow \quad \boxed{\mathbb{Q}}
		\]
		
		\item 
		\[
		\sqrt{5}+3 \in \mathbb{R}
		\quad \Rightarrow \quad \boxed{\text{Cerrada en } \mathbb{R}}
		\]
		
		\item 
		\[
		7-12=-5 \in \mathbb{Z}
		\quad \Rightarrow \quad \boxed{\text{Cerrada en } \mathbb{Z}}
		\]
		
	\end{enumerate}
	
	%------------------------------------------------
	\section*{Parte III: Propiedades}
	
	\begin{enumerate}
		
		\item 
		\[
		(17+23)+(13+7)=40+20=60
		\]
		
		\item 
		\[
		(12+18)+(8+2)=40
		\]
		
		\item 
		\[
		6(15+5)=6\cdot15+6\cdot5=120
		\]
		
		\item 
		\[
		(8-3)(20+10)=5\cdot30=150
		\]
		
		\item 
		\[
		(5+4)(14+6)=9\cdot20=180
		\]
		
		\item 
		\[
		(7+5-3)(12+8)=9\cdot20=180
		\]
		
	\end{enumerate}
	
	%------------------------------------------------
	\section*{Parte IV: Análisis de error}
	
	\[
	(14-6)\cdot3=8\cdot3=24
	\]
	
	\[
	14-(6\cdot3)=14-18=-4
	\]
	
	\[
	\boxed{24 \neq -4}
	\]
	
	Error: mala aplicación de la propiedad distributiva.
	
	Correcto:
	\[
	(14-6)\cdot3=14\cdot3-6\cdot3=24
	\]
	
	%------------------------------------------------
	\section*{Parte V: Balance hídrico}
	
	\[
	18+11-23-4=(18+11)-(23+4)=29-27=2
	\]
	
	\[
	\boxed{2 \text{ mm}}
	\]
	
	Interpretación: balance positivo.
	
	%------------------------------------------------
	\section*{Parte VI: Fertilización}
	
	\[
	150 \times 3{,}8 = 570 \text{ kg}
	\]
	
	\[
	570 \div 25 = 22{,}8 \Rightarrow \boxed{23 \text{ sacos}}
	\]
	
	%------------------------------------------------
	\section*{Parte VII: Mezclas}
	
	\[
	2{,}4 \times 5 = 12 \text{ litros de agua}
	\]
	
	\[
	\text{Total} = 2{,}4 + 12 = 14{,}4 \text{ litros}
	\]
	
	%------------------------------------------------
	\section*{Parte VIII: Costos}
	
	\[
	4(8200)+3(8200)+2(6900)
	\]
	
	\[
	=32800+24600+13800=71200
	\]
	
	\[
	\boxed{\$71.200}
	\]
	
	Forma factorizada:
	\[
	7(8200)+2(6900)
	\]
	
	%------------------------------------------------
	\section*{Parte IX: Validación}
	
	\begin{enumerate}
	
			\item \textbf{Verificación del cálculo}
			
			El técnico realiza:
			\[
			140 \times 3{,}6 = 504
			\]
			
			Verificamos:
			\[
			140 \times 3{,}6 = 140 \times \frac{36}{10} = \frac{5040}{10} = 504
			\]
			
			Por lo tanto, el cálculo es \textbf{correcto}.
			
			\item \textbf{Análisis de razonabilidad}
			
			Como estimación:
			\[
			140 \times 4 = 560
			\]
			
			Dado que $3{,}6 < 4$, el resultado debe ser menor que $560$. Como $504 < 560$, el resultado es \textbf{razonable}.
			
			\item \textbf{Análisis del resultado}
			
			\begin{itemize}
				
				\item \textbf{Unidad:}
				\[
				(\text{kg/ha}) \times (\text{ha}) = \text{kg}
				\]
				La unidad final es \textbf{kilogramos (kg)}.
				
				\item \textbf{Magnitud:}
				\[
				\frac{504}{3{,}6} = 140 \ \text{kg/ha}
				\]
				Se respeta la dosis recomendada, por lo tanto la magnitud es coherente.
				
				\item \textbf{Coherencia del contexto:}
				
				El valor obtenido representa la cantidad total de fertilizante necesaria para toda la superficie, por lo tanto es coherente con el problema.
				
			\end{itemize}
			
			\item \textbf{Cálculo de sacos}
			
			Si cada saco contiene $25 \ \text{kg}$:
			\[
			\frac{504}{25} = 20{,}16
			\]
			
			Dado que no se pueden comprar fracciones de saco, se debe redondear hacia arriba:
			\[
			20{,}16 \rightarrow 21 \ \text{sacos}
			\]
			
			Si se compraran $20$ sacos:
			\[
			20 \times 25 = 500 \ \text{kg} < 504 \ \text{kg}
			\]
			
			Esto sería insuficiente, por lo tanto:
			
			\[
			\boxed{21 \ \text{sacos}}
			\]
			
		\end{enumerate}
		
		\section*{Conclusión}
		
		El cálculo realizado por el técnico es correcto, el resultado es razonable y coherente con el contexto, y se deben comprar \textbf{21 sacos} de fertilizante, redondeando hacia arriba para asegurar la dosis completa.
		
	
\end{document}